% vim: set tw=120 fo+=t sw=2 ts=2 :
\PassOptionsToPackage{unicode=true}{hyperref} % options for packages loaded elsewhere
\PassOptionsToPackage{hyphens}{url}
%
\documentclass[12pt,letterpaper]{report}
\usepackage{version}
\excludeversion{TODO}
\usepackage{amssymb,amsmath}
\usepackage{ifxetex,ifluatex}
\usepackage[spanish]{babel}
\usepackage{fixltx2e} % provides \textsubscript
\ifnum 0\ifxetex 1\fi\ifluatex 1\fi=0 % if pdftex
  \usepackage[T1]{fontenc}
  \usepackage[utf8]{inputenc}
  \usepackage{textcomp} % provides euro and other symbols
\else % if luatex or xelatex
  \usepackage{unicode-math}
  \defaultfontfeatures{Ligatures=TeX,Scale=MatchLowercase}
\fi
% use upquote if available, for straight quotes in verbatim environments
\IfFileExists{upquote.sty}{\usepackage{upquote}}{}
% use microtype if available
\IfFileExists{microtype.sty}{%
\usepackage[]{microtype}
\UseMicrotypeSet[protrusion]{basicmath} % disable protrusion for tt fonts
}{}
\IfFileExists{parskip.sty}{%
\usepackage{parskip}
}{% else
\setlength{\parindent}{0pt}
\setlength{\parskip}{6pt plus 2pt minus 1pt}
}
\usepackage{hyperref}
\hypersetup{
            pdfborder={0 0 0},
            breaklinks=true}
\urlstyle{same}  % don't use monospace font for urls
\setlength{\emergencystretch}{3em}  % prevent overfull lines
\providecommand{\tightlist}{%
  \setlength{\itemsep}{0pt}\setlength{\parskip}{0pt}}
\setcounter{secnumdepth}{0}
% Redefines (sub)paragraphs to behave more like sections
\ifx\paragraph\undefined\else
\let\oldparagraph\paragraph
\renewcommand{\paragraph}[1]{\oldparagraph{#1}\mbox{}}
\fi
\ifx\subparagraph\undefined\else
\let\oldsubparagraph\subparagraph
\renewcommand{\subparagraph}[1]{\oldsubparagraph{#1}\mbox{}}
\fi

\usepackage{biblatex}
\addbibresource{referencias.bib}

% set default figure placement to htbp
\makeatletter
\def\fps@figure{htbp}
\makeatother

\usepackage{graphicx}
\usepackage{fancyhdr}

\pagestyle{fancy}
% \renewcommand{\footrulewidth}{0.4pt}
\fancyhf{}
\lhead{\rightmark}
\rhead{Página \textbf{\thepage}}
% \cfoot{Universidad de Chile - Facultad de Ciencias Físicas y Matemáticas\\Departamento de Ciencias de la Computación}

\frenchspacing
\title{{\small Introducción al Trabajo de Título \\  Propuesta de tema}
\\
{\Huge
Mensajes Personales en Gamificación:\\
Aplicación a la técnica Pomodoro.}
}
\author{Nicolás Salas V.\\\texttt{nikosalas@gmail.com}\\\texttt{17.478.091--9}\\Tel.: \texttt{(2) 22490790} --- \texttt{+56 9 8252 6659}}
\def\thesection{\arabic{section}}

\begin{document}

\begin{figure}[t]
\includegraphics[scale=0.83]{logo.png}
%\hspace{3.5cm}
\begin{tabular}{l}
\small Universidad de Chile\\
\small Facultad de Ciencias Físicas y Matemáticas\\
\small Departamento de Ciencias de la Computación\\
\small CC6908 Introducción al Trabajo de Título\\
\small Prof: Jérémy Barbay
\vspace{2.3cm}
\end{tabular}
\end{figure}
\maketitle

\begin{abstract} En este trabajo se busca comparar la efectividad de
incorporar mensajes personalizados (texto, audio y potencialmente
vídeo), música y otras técnicas de \emph{gamification} para mejorar la
adhesión de los usuarios a los principios de la metodología de manejo
de trabajo \emph{pomodoro}~\cite{cirillo2009pomodoro}. El resultado
del trabajo consistirá en dos partes: un producto permitiendo a un
usuario de incorporar todas las técnicas implementadas a su flujo de
trabajo, y un estudio preliminar de la efectividad de estas técnicas
separadas.
\end{abstract}

\tableofcontents
\newpage

\section{Motivación}\label{motivacion}

La metodología llamada ``\emph{Pomodoro}''~\cite{cirillo2009pomodoro}
es una técnica de manejo de trabajo que apunta al aumento de la
productividad. En una primera instancia, se trata de dividir el
trabajo en períodos de alrededor de 25 minutos para luego tomar una
pausa corta que típicamente dura 5 minutos. A esta ventana de tiempo
(30 minutos) se le asigna una \emph{unidad} llamada Pomodoro.  Luego
de cumplir cuatro pomodoros, la persona puede tomar un descanso más
largo de 25 ó 30 minutos. Finalmente, después de ese largo descanso,
se repite el ciclo. En un nivel de aplicación mas profundo, se pide al
usuario de asignar una cantidad de unidades de tiempo a cada tarea al
inicio del periodo de trabajo (sea un día, o una semana), y de revisar
la cualidad de su asignación al fin del periodo de trabajo, de maneja
a mejorar sus estimaciones futuras.

Las técnicas de \emph{gamification}\footnote{La palabra en español es
ludificación, pero para preservar el nombre que la literatura da, se
usará \emph{gamification}.}~\cite{deterding2011gamification} consisten
en la aplicación de elementos lúdicos, en general mas propios a
juegos, a contextos dichos ``de trabajo'' con un producto útil.
Desarrollada desde la década de los 80, este técnica ha encontrado
muchas aplicaciones, ya sea en Marketing~\cite{hamari2010game},
educación~\cite{desousa2014systematic-gamification},
trabajo~\cite{oravec2015gamification}, o
salud~\cite{pereira2014gamification-review}. Se espera que aplicar la
\emph{gamification} a ciertos procesos ayuda a la productividad y
motivación de la persona, pero se ha demostrado las limitaciones de
esta estrategia~\cite{hamari2014gamification}.

Si bien existen aplicaciones que han incorporado elementos de
gamification a la técnica pomodoro, hasta la fecha no se ha explorado
la incorporación de recompensas en forma de \emph{mensajes
electrónicos personalizados} (de texto, audio o video), tales como los
que han sido democratizados por aplicaciones como \texttt{WhatsApp} o
\texttt{Telegram}.

\paragraph{Hipótesis} El uso de técnicas de gamification
(tradicionales o basadas en musica o mensajes personales), ¿podría
ayudar a algunos usuarios a \emph{aprender} y \emph{seguir} la
metodología de trabajo de \emph{pomodoro}? Dentro de estas técnicas,
¿cuáles son las más potentes para la mayoría de los usuarios? ¿Cuáles
son las más potentes para cada categoría de usuarios?

\newpage
\section{Posibles soluciones}\label{posibles-soluciones}

Para estudiar el impacto de la aplicación de técnicas de gamification
sobre la adhesión y la productividad de la técnica pomodoro sobre sus
usuarios, se necesita tener una aplicación que, primero, permita la
aplicación de la técnica pomodoro y provea de mensajes útiles al
usuario y, segundo, que permita tomar estadísticas de su uso, de modo
de poder medir el impacto de la aplicación de las técnicas de
gamification introducidas.

Dentro de aplicaciones existentes que se pueden usar para pomodoro hay
bastantes~\cite{forest-stayfocused, focus-extension, pomodorium,
teamviz}. De hecho, dentro de las recién citadas algunas de ellas ya
incorporan elementos de gamification o elementos de estadística. Sin
embargo, no se ha demostrado seriamente la eficacia de la aplicación
de elementos de gamification sobre
Pomodoro~\cite{hamari2014gamification}. Este trabajo busca discernir
si dicha aplicación es positiva, negativa, o no tiene impacto alguno.

Por una parte, medir la eficiencia de pomodoro es una tarea que
aplicaciones ya realizan, aunque no de forma sistemática. Por ejemplo,
TeamViz~\cite{teamviz} muestra la cantidad de pomodoros cumplidos y
permite asociarlas con tareas, de manera de tener una estimación de
tiempo por tarea.

Por otra parte, existen aplicaciones que ya incorporan elementos de
gamification para aumentar la productividad de sus usuarios, Por
ejemplo existe \texttt{Habitica}~\cite{habitica} para propósitos
generales, \texttt{Pomodorium}~\cite{pomodorium} para Pomodoro o
\texttt{Zombies, Run!}~\cite{zombiesrun} para deportistas.

Por último, un elemento de gamification que se propone en este trabajo
es abordar la incorporación de mensajes personales (de texto, audio o
video) como recompensa por cumplir objetivos del método pomodoro. El
usuario podrá leer, ver o escuchar estos mensajes solamente si cumple
algunas metas predefinidas. Por supuesto, puede haber otros elementos
de gamification que se incorporen a esta metodología de asignación de
recompensas, pero por ahora sólo se define lo que se ha dicho
anteriormente.

Una solución que se puede usar para medir la eficiencia y el apego a
la técnica pomodoro es simplemente integrar aplicaciones que resuelvan
estos temas por separado. Escuchar algunos mensajes personalizados
desde \texttt{WhatsApp}, usar \texttt{TeamViz} para \emph{Pomodoro} y
Habitica para incorporar los elementos del juego. Si bien en papel
esto resuelve el problema planteado, es impracticable para el estudio
que se requiere. La mala interconexión entre aplicaciones generaría
cambios de contexto innecesarios que disminuirían la productividad del
usuario y la cantidad de datos proporcionados podría no ser adecuada
para estudiar el apego al método pomodoro o el aumento de
productividad.

Finalmente, para proponer la extensibilidad de un proyecto como este a
la productividad en general, se añade un módulo de música que
permitirá al usuario seleccionar listas de reproducción para momentos
de descanso o de trabajo según el período de pomodoro en que se
encuentre la persona. Esto permitirá poner algunas recompensas de
audio intercaladas en la reproducción, con la hipótesis de que ello
animará al usuario en su trabajo resultando también en un aumento de
productividad~\cite{hallam2002effects}.

\newpage
\section{Objetivos del proyecto}\label{objetivos-del-proyecto}

Este trabajo tiene un objetivo general y un desglose de varios
objetivos específicos que ayudan a cumplir ese objetivo general.

\subsection{Objetivo general}\label{objetivo-general}

Proveer una aplicación de apoyo al método Pomodoro que incorpore
mensajes personalizados de texto, audio, video y elementos del juego
(gamification) con el objetivo de aumentar la productividad de su
usuario y medir los resultados de aumento, baja o mantención de
productividad.

\subsection{Objetivos específicos}\label{objetivos-especificos}

\subsubsection{Implementación de un servicio de apoyo al método Pomodoro}\label{implementacion-de-un-servicio-de-apoyo-a-pomodoro}

El tema central de experimentación de este trabajo gira en torno a
Pomodoro. Tener una aplicación que permita a su usuario implementar el
método de forma correcta y de acuerdo a sus necesidades es esencial
para poder medir la productividad de la persona.

Dentro de los hitos de desarrollo de la aplicación que involucran este objetivo se encuentran:

\begin{enumerate}\tightlist{}
  \item \textbf{Reproductor minimal}: Puede reproducir sonido desde
carpetas indicadas por el usuario. Intercala entre archivos de las
carpetas.
  \item \textbf{Pomodoro}: Puede tomar tiempos de trabajo y de
descansos largos y cortos. Anuncia via voz y/o notificaciones los
momentos de descansos y fines de pausas largas o cortas.
  \item \textbf{Ajuste de parámetros}: Es capaz de leer parámetros de
configuración de pomodoro y personales desde un archivo de
configuración. Todavía no hay interfaz de usuario.
  \item \textbf{Interfaz de usuario minimal}: Se puede cambiar
opciones por medio de una interfaz gráfica minimal. Los cambios son
reflejados en el archivo de configuración.
  \item \textbf{Bloqueo de estación de trabajo}: La interfaz de
usuario (en adelante UI) avisa sobre hitos importantes como pausas
cortas y largas, fin de descanso. También puede bloquear la pantalla
bajo autorización del usuario.
  \item \textbf{Modificación del flujo}: Permite las opciones de
saltarse o postergar pausas por el usuario. Por supuesto esto no
bloquea la estación de trabajo.
  \item \textbf{Implementación en plataformas de escritorio}: La
aplicación cuenta con una UI funcional no necesariamente gráfica en
macOS, Linux y Windows.
\end{enumerate}

\subsubsection{Implementación de gamification sobre Pomodoro}\label{implementacion-de-gamification-sobre-pomodoro}

Con la implementación de elementos de gamification sobre Pomodoro se
puede empezar a medir la adhesión y apego al método de parte de los
usuarios. En un principio el elemento que se implementará tiene que
ver con el permiso a leer, escuchar o ver mensajes personales enviados
al usuario, pero esto no tiene por qué ser el único elemento de
gamification implementado.

Los hitos de desarrollo que definen este objetivo son:

\begin{enumerate}\tightlist{}
  \item \textbf{Grabado de comentarios}: Los usuarios pueden grabar
comentarios según la categoría del sonido. Ellas son: Anuncio de
pausas cortas, anuncio de pausas largas, anuncio de descansos, anuncio
de fines de descanso y sonidos personalizados para otros usuarios
(recompensas).
  \item \textbf{Edición de comentarios}: Los usuarios pueden editar
sus comentarios, cambiar la categoría de sus sonidos o eliminar sus
grabaciones.
  \item \textbf{Interacción entre usuarios}: Los usuarios pueden
intercambiar mensajes entre ellos según sus categorías.  Los mensajes
que tienen como propósito ser entregados como recompensa quedarán como
tal.
  \item \textbf{Implementación de recompensas}: Los usuarios reciben
sus recompensas según cantidad de pomodoros cumplidos, metas propias y
apego el método (esto es, sin atrasos ni postergaciones de pausas).
\end{enumerate}

\subsubsection{Medir resutados de productividad}\label{medir-resultados-de-productividad}

Finalmente se revisa qué tanto aumenta la productividad el uso de la
aplicación desarrollada y se comprueba la hipótesis del trabajo. El
objetivo es proveer una medición objetiva del cambio de productividad
de los usuarios y contrastarla con las opiniones de los mismos.

Los hitos que definen este objetivo son:

\begin{enumerate}\tightlist{}
  \item \textbf{Generación de estadísticas}: El programa detecta los
pomodoros cumplidos en un día y sus fracciones, la cantidad de pausas
saltadas, la cantidad de tiempo de trabajo que debió ser descanso y
atrasos en las pausas. Estos datos se guardan en un archivo para
conocimiento del usuario y para medir resultados.
  \item \textbf{Recopilación de datos}: Se recopila y se agregan todos
los datos que la aplicación guarda sobre la productividad de los
usuarios.
  \item \textbf{Presentación de resultados}: Se contrasta
estadísticamente los datos y se comprueba si la hipótesis del trabajo
es correcta o no. Esto es, si la aplicación de gamification ayuda a
los usuarios a aprender y mantener los principios de pomodoro y si la
productividad aumenta al aplicar gamification.
\end{enumerate}

\newpage
\section{Idea general de la solución}\label{idea-general-de-la-solucion}

El trabajo consiste en una aplicación que ayuda a sus usuarios a
cumplir el método de manejo de tiempo Pomodoro. La aplicación muestra
mensajes personalizados de texto, audio o video como recompensa por
apego al método. Adicionalmente, la aplicación es capaz de reproducir
listas de reproducción locales de música intercaladas con mensajes de
audio dependiendo de la etapa de pomodoro en la que se encuentre el
usuario.

En los objetivos específicos del proyecto se introdujo muy
sucintamente los hitos del trabajo a desarrollar. En esta sección
primero se profundiza en dichas definiciones y luego se resuelven
posibles problemas que puedan aparecer durante el desarrollo del
trabajo.

\subsection{Hitos}

\subsubsection{Reproductor minimal}

El producir un reproductor de audio minimal permitirá a la aplicación
incorporar el elemento de gamification propuesto en forma de audio,
anunciar eventos relacionados a pomodoro como pausas largas y cortas y
descansos y también permitirá al usuario escuchar música local
intercalada con eventos importantes de pomodoro si así lo desea.

\subsubsection{Pomodoro}

El elemento principal de Pomodoro es el manejo del tiempo. En este
hito se implementan contadores de tiempo que cuando se cumplen generan
un evento notifica a su usuario que el evento ha cumplido su
tiempo. Las notificaciones son de escritorio o en modo de audio.

\subsubsection{Ajuste de parámetros}

Luego de tener pomodoro y reproducción de audio, la incorporación de
parámetros permitirá a sus usuarios ajustar, por ejemplo, los tiempos
de descanso y de trabajos. Esto se hace porque el método pomodoro
especifica que los tiempos de atención no son iguales para todas las
personas. Los parámetros tentativos son:

\begin{itemize}\tightlist{}
  \item Tiempo de enfoque en trabajo
  \item Tiempos de pausas cortas
  \item Tiempo de descanso largo
  \item Carpeta de música de trabajo
  \item Carpeta de música de pausa
  \item Carpeta con mensajes personalizados
\end{itemize}

\subsubsection{Interfaz de usuario minimal}

Posteriormente se empieza a elaborar una interfaz de usuario que en
este hito permite a sus usuarios ajustar los parámetros definidos
anteriormente según su volición. Estos cambios se reflejan en el
archivo de configuración que la aplicación provee.

\subsubsection{Bloqueo de estación de trabajo}

Para ayudar a cumplir con la técnica pomodoro, es útil que la estación
de trabajo se bloquee durante los períodos de pausa. Según el autor
del método, es mejor salir de la estación de trabajo durante las
pausas para descansar la mente y los ojos.

\subsubsection{Modificación del flujo de trabajo}

El usuario podría querer postergar pausas o simplemente saltárselas
por razones de tiempo o concentración. Esta aplicación permite que el
usuario realice estas tareas aunque no sea apropiado para el
método. El uso de esta característica de la aplicación después sirve
para medir el apego a pomodoro.

\subsubsection{Generación de estadísticas}

La aplicación grabará datos estadísticos de uso y aplicación de la
técnica según lo que el usuario decida hacer con la ejecución del
método. Se mide la cantidad de postergaciones de descanso, los atrasos
en los descansos y cantidad de pomodoros cumplidos por día. Estos
datos servirán en el futuro para que el usuario mismo vea cómo se ha
apegado a la técnica y para medir resultados.

\subsubsection{Incorporación de gamification}

Los usuarios podrán intercambiar mensajes dirigidos de texto, audio o
video que serán usados como recompensa por cumplir objetivos (por
ejemplo, hacer cuatro pomodoros). A cada mensaje se asigna una
categoría y cada usuario puede grabar nuevos mensajes que se pueden
compartir entre usuarios. Los mensajes se muestran dependiendo de la
categoría del mismo.  Algunas categorías son: Anuncio de pausas largas
y cortas, vuelta a la estación de trabajo y recompensas.

\subsubsection{Implementación en plataformas de escritorio}

Con el objetivo de hacer esta aplicación ubicua, será lanzada para
todas las plataformas de escritorio populares: Windows, macOS y
Linux. La interfaz de usuario será minimal pero ampliamente funcional.


\subsubsection{Grabado de comentarios}

El grabado de comentarios en una característica fundamental del
programa. Tenerla permite que se empiece a implementar la
gamification, pues a medida que avance el desarrollo los comentarios
grabados se pueden reproducir o compartir con otros usuarios de la
aplicación. Las categorías de grabado son: Anuncio de pausas cortas,
anuncio de pausas largas, anuncio de descansos, anuncio de fines de
descanso y sonidos personalizados para otros usuarios (recompensas).

\subsubsection{Edición de comentarios}

Los comentarios tienen ciertos \emph{tags}, que fueron enunciados
anteriormente. Un usuario puede equivocarse al categorizar una
grabación o puede querer simplemente eliminar sus grabaciones.

\subsubsection{Interacción entre usuarios}

Los usuarios pueden intercambiar mensajes entre ellos según sus
categorías. Los mensajes que tienen como propósito ser entregados como
recompensa quedarán como tal. El intercambio es otro paso fundamental
de gamification. Con él los usuarios pueden motivar a otras usuarios
con mensajes previamente grabados.

\subsubsection{Implementación de recompensas}

Lo último que quedaría es la implementación de \textbf{la entrega de
recompensas} a los usuarios según reglas: por ejemplo, pomodoros
cumplidos, disminución de atrasos o metas propias de los usuarios. Con
este hito implementado los usuarios finalmente podrán experimentar la
gamification en pomodoro.

\subsection{Soluciones a posibles riesgos del proyecto}

Para disminuir el riesgo del proyecto, a continuación se introduce
posibles soluciones a problemas que pueden, en primera instancia, no
ser triviales de resolver.


\subsubsection{Multiplataforma}\label{multiplataforma}

La multiplataforma es un tema fundamental para este proyecto. Mientras
más multiplataforma sea el proyecto, más fácil será evaluar
resultados, pues más usuarios tendrán acceso a la aplicación
desarrollada.

Para disminuir el riesgo de este problema, se ocupará un lenguaje que
provea como característica el ser multiplataforma.  Dos candidatos que
resuelven este problema son Java (via su librería
\texttt{swing}~\cite{javaswing}) y \texttt{Node.js}~\cite{nodejs} vía
sus librerías \texttt{electron}~\cite{electron} y \texttt{React\
Native}~\cite{react-native}.

La solución es simplemente usar \texttt{Node.js}. El que \texttt{React
Native} y \texttt{electron} provean una arquitectura en la que se
construye interfaces de usuario con HTML facilita enormemente la tarea
de exportar el código a ambientes móviles (iOS/Android) mientras que
\texttt{Electron} facilita la multiplataforma en aplicaciones de
escritorio (macOS/Linux/Windows).

\subsubsection{Bloqueo de estación de trabajo}\label{bloqueo-de-estacion-de-trabajo}

Podría parecer que este es un problema que no se resuelve
fácilmente. Workrave~\cite{workrave} es un programa licenciado bajo la
GNU General Public License v3 que resuelve este problema. Tomar la
parte de este programa que resuelve este problema disminuye la
incertidumbre sobre cómo resolver este problema, pues, aunque pueda
haber otras soluciones, ya existe una forma de abordar el problema
\emph{a priori}.

\newpage
\section{Metodología}\label{metodologuxeda}

En este apartado se describe el trabajo a realizar para cumplir los
objetivos

\subsection{Iterar los ciclos de desarrollo}\label{iterar-los-ciclos-de-desarrollo}

Anteriormente se definieron los hitos que marcan el proceso de
desarrollo de esta aplicación. En este proceso cada hito supera en
funcionalidad al anterior y por sobre todo, \textbf{es usable}. Antes
de cada iteración se debe estimar el tiempo que la implementación
tomará, de modo de ajustar las horas de trabajo de la aplicación para
cumplir el objetivo final.

\subsection{Medición de resultados}\label{mediciuxf3n-de-resultados-1}

Con al menos la plataforma de escritorio de la aplicación, se prueba
la aplicación con personas que tengan interés en aplicar la técnica
pomodoro de manejo de tiempo o bien estén interesadas en aplicarle
elementos de gamification a su proceso productivo. En este período, se
miden los resultados completos en una semana completa de trabajo, dos
veces por separado. Se propone esta metodología para ocupar la primera
semana como medida base de diagnóstico y acostumbramiento y la segunda
semana como medición de mejora sobre la técnica. Los resultados a
medir son:

\begin{itemize}
  \item Cantidad de pomodoros cumplidos
  \item Tiempo de atraso en pausas
  \item Uso de elementos de gamification provistos
  \item Opinión del usuario sobre el uso de Pomodoro
\end{itemize}

El procesamiento de estos resultados da la posibilidad de estudiar la
evolución en el tiempo de los usuarios en la adhesión y el apego a
pomodoro. El contraste de esas medidas con la opinión de cada usuario
permitirán dar un resultado general de cómo cambia la productividad
del usuario.

Para ver cómo impacta la gamification en todo este proceso, durante
las fases de medición se comenzará sin la característica de
gamification y se introducirá durante la segunda fase de pruebas.

En términos de la muestra escogida y por asuntos de tiempo, se propone
una muestra de entre 20 a 30 personas, deseablemente de todas las
edades, pero -de nuevo-, por asuntos de tiempo, el rango etáreo es
solamente deseable que sea uniforme.

\subsection{Planteamiento de posibles extensiones}\label{planteamiento-de-posibles-extensiones}

Dependiendo del tiempo que todo lo anterior haya tomado, puede
implementarse algunas de las extensiones propuestas en el apéndice o
exportarse a alguna otra plataforma. Si esto no llegase a ser posible,
se especificará debidamente cómo se puede extender esta aplicación
para otros contextos.

\newpage
\section{Apéndice}\label{sec:apendice}

\subsection{Posibles extensiones}\label{subsec:posibles-extensiones}

Pomodoro no es el único contexto en el que este trabajo puede ser
útil, en este trabajo se ha escogido como la técnica para medir la
efectividad de aplicar elementos de gamification a esta. A
continuación se mencionan algunos ejemplos en los que esto puede
también ser útil.

\subsubsection{Deportes}

Existen algunas aplicaciones ya mencionadas~\cite{zombiesrun}, que
incorporan elementos de gamification en deportes, pero se ha dejado de
lado el factor humano y la influencia de mensajes de seres queridos
sobre la persona que utiliza estas aplicaciones. Un tema de estudio
relacionado es cómo afecta una aplicación como la que este trabajo
describe sobre deportistas y si efectivamente los ayuda a mejorar su
rendimiento.

\subsubsection{Hospitales}

Las personas en situación hospitalaria también pueden ser objeto de
estudio de una aplicación como ésta, sobre todo en pacientes que de
alguna forma pueden sentirse solos. Proveer un mecanismo de
gamification combinado con mensajes personales podría mostrar un
aumento en el ánimo.

\newpage
\section{Referencias}
\printbibliography[heading=none]

\end{document}

%%% Local Variables:
%%% mode: latex
%%% TeX-master: t
%%% End:
